\documentclass[aspectratio=169,14pt]{beamer}
\usetheme{metropolis}
\usepackage{booktabs}
\usepackage{xcolor}
\usepackage{graphicx}
\usepackage{hyperref}
\usepackage{pgfpages}

% Figures live with the notebook that produces them (single source, 300 dpi).
\graphicspath{{../eda-5/plots/}}

% Slide number (current / total) in the footer
\setbeamertemplate{frame numbering}[fraction]

% --- Out-of-bounds detection --------------------------------------------
% Build with  -usepretex='\def\strict{1}'  (see build.sh --strict) to draw a
% black rule wherever content overflows the slide, so an out-of-bounds box is
% impossible to miss in the PDF. Off by default, so the shipped deck is clean.
\ifdefined\strict
  \overfullrule=8pt
  \hbadness=1000 \vbadness=1000
\fi
% ------------------------------------------------------------------------

% --- Speaker notes -------------------------------------------------------
% Hidden by default, so the handout PDF the audience reads carries no notes.
% To get the presenter version (slide on the left, the data notes on the
% right), comment the line below and uncomment the second one, then rebuild.
\setbeameroption{hide notes}
% \setbeameroption{show notes on second screen=right}
% -------------------------------------------------------------------------

% Every plot is included with the SAME box and keepaspectratio, so no figure
% is ever stretched or squashed and none is resized differently from the rest.
\newcommand{\plotframe}[4]{%
  \begin{frame}{#1}
    \begin{center}
      \includegraphics[width=\linewidth,height=0.58\textheight,keepaspectratio]{#2}
    \end{center}
    \vspace{-0.5em}\footnotesize #3
  \end{frame}%
  \note{#4}}

% Taller box for the dense per-community network plots, which need the height
% to stay readable; one-line caption keeps it inside the slide.
\newcommand{\plotframeT}[4]{%
  \begin{frame}{#1}
    \begin{center}
      \includegraphics[width=\linewidth,height=0.66\textheight,keepaspectratio]{#2}
    \end{center}
    \vspace{-0.4em}\footnotesize #3
  \end{frame}%
  \note{#4}}

\title{Enron Email Corpus -- Week 5}
\subtitle{Cleaning the corpus, reworking the review}
\author{Data Mining Lab SoSe 26 -- Task 5}
\institute{TUM, Department of Computer Science}
\date{June 3, 2026}

\begin{document}

% metropolis' title block packs ~20pt taller than \textheight, but the title
% frame has room below it so nothing clips. Relax the overflow tolerance for
% this one frame only, so the detector flags real out-of-bounds boxes and not
% this known-benign one. Every content frame keeps the strict default.
{\vfuzz=24pt \vbadness=\maxdimen \maketitle}
\note{%
  Plan: about 15 minutes, 13 content slides, roughly a minute each. Two
  halves. First half: we cleaned the whole corpus this week, and the cleaning
  by itself overturns two numbers from the week-2 report. Second half: with
  clean data we redo the review points (the 24-hour plot, the social graph on
  people, the word charts on one scale, deletion) and add the one thing we
  never had before, conversation threads. Roles: plots by Max, report and this
  deck by Len.}

% ---------------------------------------------------------------
\begin{frame}{From raw files to a clean corpus}
  \small
  Last week's raw-file problems: about half were duplicates, bodies carried
  quoted history, addresses were split per person, and there were no thread
  headers. We cleaned the corpus in four passes into \texttt{eda-5/clean/}:
  \begin{itemize}
    \item \textbf{Deduplicate} to one row per distinct message
    \item \textbf{Clean bodies}: strip quoted history, HTML, disclaimers
    \item \textbf{Resolve addresses}: spelling variants to one person
    \item \textbf{Reconstruct threads}: subject + participant, 30-day gap
  \end{itemize}
  Result: \textbf{252{,}022} messages and \textbf{1.37\,M} recipient rows.
\end{frame}
\note{%
  The framing: none of last week's plots were wrong arithmetic, they were
  honest counts of the wrong unit, the file instead of the message. So the fix
  is to clean once and rerun, not patch each plot. Dedup key: same sender,
  subject, exact send time and recipient counts; checked against body hashes,
  40 of 40 identical. The body cleaner cuts only on header-anchored reply
  markers; an earlier version that also cut bare divider lines destroyed about
  8,900 real newsletters, so we pulled that back. Threads: the release has no
  In-Reply-To / References, so we rebuild from a normalised subject plus an
  overlapping participant, capped at a 30-day gap so unrelated ``fyi'' notes do
  not chain across years. Full procedure and limits are in DATASET\_CLEANING.md.}

% ---------------------------------------------------------------
\plotframe{Half the files are duplicates}{b1_dedup_waterfall.png}{%
  \textbf{517{,}401} files collapse to \textbf{252{,}022} distinct messages
  (\textbf{51\%} duplication). Every later count is on the distinct messages.}{%
  About 203,000 of the removed copies are within one mailbox, the same message
  filed in both inbox and a project folder, a normal Outlook habit. The other
  roughly 63,000 are one message sitting in several people's mailboxes because
  they were all on it. So a raw file count inflates broadcasts and carefully
  filed memos many times over.}

% ---------------------------------------------------------------
\plotframe{What dedup changes}{b1b_dedup_corrections.png}{%
  Two week-2 numbers were artifacts. \texttt{topic/project} folders drop from
  \textbf{37\%} to about \textbf{11\%}; the largest mailbox flips from
  \texttt{kaminski-v} to \texttt{dasovich-j}. Left bars are \emph{share} of the
  corpus, not counts.}{%
  Project and topic folders looked huge because they are where people re-file
  copies of mail they already hold in inbox or sent, so they carry more than
  their share of duplicates. Once each message counts once, sent and inbox are
  the real bulk. kaminski-v was the raw ``largest mailbox'' at 28,465 files,
  but it duplicated more heavily than dasovich-j, so the ranking was an
  artifact of duplication, not size.}

% ---------------------------------------------------------------
\plotframe{How complete is the data}{b2_metadata_coverage.png}{%
  Identity and timing are near-complete: sender \textbf{100\%}, timestamp
  \textbf{99.9\%}, subject \textbf{96\%}. The threading headers are gone.}{%
  Message-ID reads as 100 percent present but the values are machine-generated
  placeholders with no link to the message they answer, and In-Reply-To and
  References were stripped before release. That is the single biggest gap: the
  corpus knows every message but not which message answered which. It is why
  threads have to be reconstructed rather than read off.}

% ---------------------------------------------------------------
\plotframe{The 24-hour activity plot, fixed}{a1_weekday_hour_dual.png}{%
  Top row (blue): one shared count scale, honest about volume. Bottom row
  (green): each year on its own scale, honest about the daily shape. Empty
  cells are white.}{%
  The old plot put all four years on one dark-for-zero colour bar, so quiet
  years read as empty and empty grids read as full. Splitting it fixes both.
  The daily block is the same every year, about 14:00 to 23:00 UTC, which is
  mid-morning to late afternoon in Houston. Weekends are white. The one
  exception is 2002, thinner and shifted, a company being wound down.}

% ---------------------------------------------------------------
\plotframe{A keyhole onto Enron}{b4_recipient_coverage.png}{%
  The 150 mailboxes wrote to about \textbf{33{,}000} distinct internal
  accounts, but only about \textbf{17\%} of those have a mailbox of their own
  in the release.}{%
  This is the standing caveat behind every network and rate figure. The graph
  is dense and complete among the 150 captured people and increasingly partial
  as it reaches everyone else. We see one side of most conversations. The 150
  were released because they were central to the FERC investigation, so they
  are far more connected to each other than a random 150 employees would be.}

% ---------------------------------------------------------------
\plotframe{Addresses are not people}{c1_addresses_per_person.png}{%
  Of 150 owners, \textbf{104} write from one address, \textbf{31} from two,
  \textbf{6} from three or four. The multi-address cases are messy formatting
  of one name, not extra accounts.}{%
  Heaviest case is Vince Kaminski: kaminski, j.kaminski and j..kaminski.
  One clarification on the feedback wording: here a mailbox really does equal a
  person, the only repeated surnames (mckay, ring, whalley) are each two
  different people. So nothing merges at the folder level; the multiplicity is
  entirely in the addresses, and the next plots collapse them so each person is
  one node.}

% ---------------------------------------------------------------
\plotframe{The reciprocal graph, grouped}{d2_reciprocal_cluster.png}{%
  \textbf{150} owners, \textbf{4{,}768} two-way links, coloured by a modularity
  split. Dense, but not structureless: five groups separate out.}{%
  A link here means the two people mail each other in both directions. We lay
  each community out as its own blob and drop the small executive and
  government-affairs core in the centre, where its cross-group links are
  shortest, so the colours group instead of scattering. The density is the
  headline, but remember it is partly the FERC selection, not a natural office.}

% ---------------------------------------------------------------
\plotframeT{Working group 1: trading floor + research}{d3_1.png}{%
  The largest community (56): the trading desk, research and risk around it.}{%
  Kaminski, Buy, Arnold, Delainey, Beck. This is the core book-running floor;
  research and risk sit inside it rather than beside it, which is why they
  cluster together.}

% ---------------------------------------------------------------
\plotframeT{Working group 2: West-power / real-time desk}{d3_2.png}{%
  36 people: the West-power / real-time desk that ran the California book.}{%
  Forney, Presto, Campbell. This is the desk behind the California content in
  the next section; it is tightly internally connected and links up to the
  executive core.}

% ---------------------------------------------------------------
\plotframeT{Working group 3: ENA legal + East gas}{d3_3.png}{%
  29 people: the ENA legal department and the East gas desk it served.}{%
  Jones, Shackleton, Haedicke, Mann, Germany, Nemec. Legal and its client desk
  show up as one community because they corresponded constantly.}

% ---------------------------------------------------------------
\plotframeT{Working group 4: pipelines / utility ops}{d3_4.png}{%
  18 people: the pipelines and utility-operations group.}{%
  Lokay, Hayslett, McConnell, Scott, Watson. The most self-contained group, the
  operational side rather than trading or legal.}

% ---------------------------------------------------------------
\plotframeT{Working group 5: executives + government affairs}{d3_5.png}{%
  The smallest (11), and the core that ties the others together.}{%
  Lay, Skilling, Whalley with Dasovich, Kean, Shapiro, Steffes and general
  counsel Derrick. Small but central; this is the group sitting in the middle of
  the full graph two slides back.}

% ---------------------------------------------------------------
\plotframe{California vs EES, on one scale}{e1_california_ees_samescale.png}{%
  Now on the full corpus, not a 25k sample: about \textbf{12{,}500} bodies
  mention California against about \textbf{3{,}900} for EES. Same axis, so the
  gap is real.}{%
  California is the far heavier topic, and its companion words are the crisis
  vocabulary, power, energy, state, market, price. EES is shorter and its words
  are internal, organisational, the retail-services unit. Because the bodies
  are cleaned, the words are the writers' own and not quoted history, and
  flagged signatures and disclaimers are excluded so they cannot crowd the top.}

% ---------------------------------------------------------------
\plotframe{What gets deleted}{f2_deletion_topic_preference.png}{%
  Deletion is topic-driven. Base rate about \textbf{14\%}. Far above it sit one
  protest campaign (\texttt{donate}, \texttt{stock}, \texttt{lay}) and machine
  noise; business words are kept.}{%
  The over-deleted words belong to a mass campaign demanding Ken Lay donate his
  stock-sale proceeds, thousands of near-identical messages swept straight into
  his deleted folder, plus automated traffic like headlines and over-quota
  warnings. Under-deleted are deal words: nda, raptor, abb, allegheny. On
  timing (the other plot, f1), deletion just tracks overall volume and peaks in
  late 2001, with no clean jump at any single collapse date. We also tested
  whether leaving drove deletion: 76\% of mailboxes are end-loaded, median ratio
  1.65, but the deleted-items buffer is emptied periodically, which produces
  that pattern on its own, so it is not conclusive.}

% ---------------------------------------------------------------
\plotframe{Conversations, finally}{h1_threads.png}{%
  Reconstruction gives \textbf{40{,}306} threads. Replies are fast: median
  about \textbf{3 hours}, three quarters inside a day. Sizes are heavily
  skewed, mostly short.}{%
  A quarter of replies land within twenty minutes. The long tail is slow-burn
  threads that resume after a pause; because our heuristic caps gaps at 30
  days, one of those gets split in two, but that barely affects reply-speed or
  size description. The single thread over a thousand messages is the protest
  campaign from the previous slide, a blast rather than a conversation.}

% ---------------------------------------------------------------
\begin{frame}{Where we stand}
  The corpus is now deduplicated, body-cleaned, address-resolved and
  thread-aware: \textbf{252{,}022} messages from 150 mailboxes, 1999--2002.

  \medskip
  Three week-2 numbers moved once we counted messages, not files: the folder
  mix inverts, the largest mailbox re-ranks, and the social core is denser than
  it looked. The new material is threads, fast replies and a heavy-tailed size
  distribution.

  \medskip
  \small The standing caveat is coverage, and the reconstructed threads are a
  heuristic standing in for headers that no longer exist.
\end{frame}
\note{%
  If asked what is next: this closes the descriptive phase. With a clean,
  thread-aware table we can move to predictive work, the schedule turns there
  next. If pressed on the cleaning's weakest point, say it is thread
  reconstruction, because the headers are gone for good and we are inferring
  them.}

% ---------------------------------------------------------------
\begin{frame}{}
  \begin{center}\Large Questions?\end{center}
  \vfill
  \small\centering Full write-up: \texttt{eda-5/DATASET\_CLEANING.md} and
  \texttt{eda-5/enron\_eda\_week4.ipynb}
\end{frame}

\end{document}
